\section{BRS symmetry}\label{sec:brs}

We now proceed with the general discussion that will lead to
the proof of finiteness of the correlation functions of the bulk
fields to all orders in the gauge coupling.
As a first step,
the BRS symmetry of the theory in $D+1$ dimensions
is reviewed in this section.

The gauge fixing discussed in sect.~\ref{ssec:gaugefixing}
follows the standard procedure and therefore leads to
a theory with a BRS symmetry.
This symmetry acts on the boundary fields in the usual way, but
the transformation of the bulk fields requires the solution of
the diffusion equation \eqref{eq:3.2} and consequently tends to be non-local.
In their work on the renormalization
of the Langevin equation, Zinn--Justin and Zwanziger \cite{ZinnZwanziger}
however showed that the locality of the transformation can be
restored by introducing additional ghost fields.

\subsection{BRS transformation of the boundary fields}

The BRS variation of the unrenormalized
fields $A_{\mu}$, $\cbar$ and $c$ is defined by \cite{BRSI,BRSII}
\begin{align}
   \dbrs A_{\mu}=D_{\mu}c,
   \label{eq:6.1}\\[2pt]
   \dbrs c=-c^2,
   \label{eq:6.2}\\[2pt]
   \dbrs\cbar=\lambda_0\partial_{\mu}A_{\mu}.
   \label{eq:6.3}
\end{align}
Note that the components $c^a$ and $\cbar^a$ of the ghost fields
and the operator $\dbrs$
anti-commute with one another.
The product on the right of eq.~\eqref{eq:6.2}, for example,
is given by
\begin{equation}
   c^2=c^ac^bT^aT^b=\frac{1}{2}c^ac^bf^{abc}T^c,
   \label{eq:6.4}
\end{equation}
and the Leibniz rule for the operator $\dbrs$
must take its anti-commuting character
into account.
The action $S+\Sgf+\Sgh$
and the measure in the functional integral
of the theory in $D$ dimensions are then
easily shown to be BRS invariant.

\subsection{Bulk ghost fields}

The additional ghost fields $d(t,x)$ and $\dbar(t,x)$
mentioned above live in $D+1$ dimensions but are otherwise
of the same kind as the Faddeev--Popov ghosts.
Their action is
\begin{equation}
   \Sflgh=-2\int_{0}^{\infty}\rmd t\int\rmd^Dx\,
   \tr\bigl\{
   \dbar(t,x)\bigl(\partial_td-\alpha_0 D_{\mu}\partial_{\mu}d\bigr)(t,x)
   \bigr\}
   \label{eq:6.5}
\end{equation}
and the $d$ field is required to satisfy the boundary condition
\begin{equation}
   \left.d\right|_{t=0}=c.
   \label{eq:6.6}
\end{equation}
No boundary condition is imposed on the $\dbar$ field, which, in many respects,
plays a r\^ole similar to the Lagrange-multiplier field $L_{\mu}$
in the case of the gauge field.

The propagator of the ghost fields,
\begin{align}
   \langle\dtilde^a(t,p)\dbartilde^b(s,q)\rangle=
   (2\pi)^D\delta(p+q)\delta^{ab}\theta(t-s)\tilde{K}_{t-s}(p)
   +\rmO(g_0^2),
   \label{eq:6.7}\\[3pt]
   \tilde{K}_{t}(p)=\rme^{-\alpha_0tp^2},
   \label{eq:6.8}
\end{align}
can be worked out following the steps taken in sect.~\ref{ssec:propagators4}.
Graphically the propagator is represented by a dashed line,
\begin{equation}
  \raisebox{-0.55cm}{\includegraphics[width=1.9cm]{images/prop4}}%
  \;=\;\delta^{ab}\theta(t-s)\tilde{K}_{t}(p),
  \label{eq:6.9}
\end{equation}
where the arrow is drawn in the direction from $\dbar$ to $d$.
There is a mixed propagator,
\begin{align}
   \langle\dtilde^a(t,p)\cbartilde^b(q)\rangle=
   (2\pi)^D\delta(p+q)\delta^{ab}g_0^2\Dtilde_t(p)
   +\rmO(g_0^4),
   \label{eq:6.10}\\[3pt]
   \Dtilde_t(p)={1\over p^2}\rme^{-\alpha_0tp^2},
   \label{eq:6.11}
\end{align}
as well, which coincides with the $c\cbar$ propagator at $t=0$ and
is therefore represented by the same graphical symbol
(a directed dotted line).
The $c\dbar$ two-point function, on the other hand,
vanishes to all orders.

The action
\begin{align}
   \left.\Sflgh\right|_{\rm interaction}=
   -\int_0^\infty\rmd t\int_{p,q,r}(2\pi)^D\delta(p+q+r)
   \notag\\[2pt]
   \times\vbulk{1,1}(p,q,r)^{abc}_{\mu}
   \tilde{B}^a_{\mu}(t,-p)\dbartilde(t,-q)^b\dtilde(t,-r)^c,
   \label{eq:6.12}
\end{align}
also gives rise to a new flow vertex $\vbulk{1,1}$.
In the Feynman diagrams,
it is represented by an open circle as the other flow vertices.
The explicit expression for the vertex is given in appendix~\ref{app:B}.

As in the case of the flow-line loops discussed in
sect.~\ref{ssec:flowlines}, it
is possible to show that diagrams with $d\dbar$ loops vanish.
Note that ghost loops with
mixed propagators do not exist, because there is no $c\dbar$ propagator
and the vertices only couple
$c$ to $\cbar$ or $d$ to $\dbar$. The only non-zero
ghost loops are thus the
usual ones of the theory at flow time zero.

\subsection{BRS variation of the fields in the bulk}

The BRS symmetry acts on the bulk fields
according to \cite{ZinnZwanziger}
\begin{align}
   \dbrs B_{\mu}=D_{\mu}d,
   \label{eq:6.13}\\[2pt]
   \dbrs L_{\mu}=[L_{\mu},d],
   \label{eq:6.14}\\[2pt]
   \dbrs d=-d^2,
   \label{eq:6.15}\\[2pt]
   \dbrs\dbar=D_{\mu}L_{\mu}-\{d,\dbar\}.
   \label{eq:6.16}
\end{align}
Note that $\dbrs B_{\mu}=\dbrs A_{\mu}$ and
$\dbrs d=\dbrs c$ at the boundary, as must be the case
in view of the boundary conditions \eqref{eq:2.3} and \eqref{eq:6.6}.

In order to show that the BRS variation of the
bulk action $\Sfl+\Sflgh$ vanishes, it is helpful to
introduce the fields
\begin{align}
   E_{\mu}=\partial_tB_{\mu}-D_{\nu}G_{\nu\mu}
   -\alpha_0D_{\mu}\partial_{\nu}B_{\nu},
   \label{eq:6.17}\\[3pt]
   e=\partial_td-\alpha_0 D_{\mu}\partial_{\mu}d.
   \label{eq:6.18}
\end{align}
After some algebra, one finds that
\begin{align}
   \dbrs E_{\mu}=[E_{\mu},d]+D_{\mu}e,
   \label{eq:6.19}\\[3pt]
   \dbrs e=-\{e,d\},
   \label{eq:6.20}
\end{align}
and the invariance of the bulk action is then easily established.

\begin{figure}[t]
\centerline{\includegraphics[width=10.0cm]{images/diagram9}}
\caption{Tree diagrams contributing to the correlation function on the left
of eq.~\eqref{eq:6.23} (diagrams 1 and 2) and to the one on the right
of the equation (diagrams 3,4 and 5).
The momenta are ingoing and satisfy $p+q+r=0$. In diagram $3$,
the momentum $p$ flows into a vertex representing
the insertion of the field $[B_{\mu},d]$, while in the case
of the diagrams $4$ and $5$ the external lines with momentum $p$
are multiplied by $ip_{\mu}$
(thus representing the field $\partial_{\mu}d$).
}
\label{fig:brsdiag}
\end{figure}

\subsection{Examples of BRS identities}

Since the functional integration measure is also invariant, it follows
that $\langle\dbrs\mathcal{O}\rangle=0$ for any product
$\mathcal{O}$ of the basic fields.
This leads to identities such as
\begin{align}
   \lambda_0\langle
   \partial_{\mu}A_{\mu}^a(x)\partial_{\nu}A_{\nu}^b(y)
   \rangle=g_0^2\delta^{ab}\delta(x-y),
   \label{eq:6.21}\\[3pt]
   \lambda_0\langle B^a_{\mu}(t,x)\partial_{\nu}A^b_{\nu}(y)\rangle=
   -\langle(D_{\mu}d)^a(t,x)\cbar^b(y)\rangle,
   \label{eq:6.22}\\[3pt]
   \lambda_0\langle
   B^a_{\mu}(t,x)\partial_{\nu}A^b_{\nu}(y)\partial_{\rho}A^c_{\rho}(z)
   \rangle=
   -\langle(D_{\mu}d)^a(t,x)\cbar^b(y)\partial_{\rho}A^c_{\rho}(z)\rangle,
   \label{eq:6.23}
\end{align}
the first of them being the familiar gauge Ward identity in the
theory in $D$ dimensions,
which determines the longitudinal part of the gauge-field
two-point function.

It is instructive to check eqs.~\eqref{eq:6.22} and \eqref{eq:6.23}
at tree level of perturbation theory. The first
equation relates the longitudinal part of the mixed
gauge-field propagator to the mixed ghost propagator, an identity
that is immediate from the expressions for the propagators.
The other equation involves the $3$-point vertices
(see fig.~\ref{fig:brsdiag}) and a non-trivial
cancellation among various contributions (appendix~\ref{app:C}).
